% !TEX root = ../main.tex

\chapter{Deskripsi Persoalan}

\section{Latar Belakang}
Tugas besar ini membahas implementasi \textit{Convolutional Neural Network} (CNN), \textit{Simple Recurrent Neural Network} (RNN), dan \textit{Long Short-Term Memory} (LSTM) untuk dua konteks yang berbeda: klasifikasi citra dan \textit{image captioning}. Berbeda dengan pendekatan \textit{end-to-end} menggunakan framework, pada tugas ini komponen \textit{forward propagation} utama diimplementasikan \textit{from scratch} menggunakan NumPy, kemudian diverifikasi menggunakan bobot hasil pelatihan model Keras.

Pendekatan ini dipilih agar proses inferensi tidak diperlakukan sebagai \textit{black box}. Untuk CNN, implementasi menekankan perbedaan antara parameter \textit{shared} (\texttt{Conv2D}) dan non-\textit{shared} (\texttt{LocallyConnected2D}). Untuk captioning, implementasi menekankan perilaku urutan waktu pada decoder RNN/LSTM dengan skema \textit{pre-inject} (fitur gambar diinjeksi sebagai $x_{-1}$) sesuai \textit{Show and Tell}.

Selain implementasi inti, proyek ini juga mengevaluasi beberapa komponen bonus yang relevan secara praktis, yaitu \textit{batch inference}, \textit{backward propagation} layer tertentu, \textit{beam search decoder}, \textit{init-inject} captioning, serta visualisasi \textit{Grad-CAM}.

\section{Tujuan}
Secara umum, tugas ini bertujuan mengembangkan pipeline \textit{deep learning} yang modular dan dapat dianalisis secara transparan. Secara spesifik, tujuan yang ingin dicapai adalah:

\begin{enumerate}
    \item Mengimplementasikan layer CNN, RNN, dan LSTM \textit{from scratch} untuk kebutuhan inferensi serta kompatibel dengan bobot Keras melalui metode \texttt{set\_weights()}.
    \item Membandingkan performa \textit{Keras vs from scratch} pada tugas klasifikasi citra (macro F1) dan captioning (BLEU-4, METEOR, waktu eksekusi).
    \item Menganalisis pengaruh hyperparameter utama terhadap kinerja model pada kedua domain tugas.
    \item Mengevaluasi dampak variasi arsitektur dan decoding strategy (\textit{pre-inject/init-inject}, \textit{greedy/beam search}) terhadap kualitas caption.
\end{enumerate}

\section{Dataset dan Ruang Lingkup}
Proyek menggunakan dua dataset yang berbeda sesuai spesifikasi tugas.

\subsection{Intel Image Classification (CNN)}
Dataset Intel Image Classification berisi sekitar 25.000 gambar dengan 6 kelas:
\textit{buildings, forest, glacier, mountain, sea, street}. Split \textit{train}, \textit{validation}, dan \textit{test} menggunakan pembagian yang telah disediakan dataset.

\begin{table}[H]
    \centering
    \caption{Ringkasan Task CNN}
    \begin{tabular}{ll}
        \toprule
        \textbf{Komponen} & \textbf{Spesifikasi} \\
        \midrule
        Task & Klasifikasi citra multi-kelas (6 kelas) \\
        Arsitektur utama & Conv2D / LocallyConnected2D + Pooling + Head + Dense \\
        Loss & Sparse Categorical Crossentropy \\
        Optimizer (Keras) & Adam \\
        Metrik evaluasi utama & Macro F1-score \\
        \bottomrule
    \end{tabular}
    \label{tab:ringkasan-cnn}
\end{table}

\subsection{Flickr8k (RNN/LSTM Captioning)}
Dataset Flickr8k berisi 8.092 gambar, masing-masing dengan 5 caption. Pembagian data yang digunakan: 6.000 \textit{train}, 1.000 \textit{validation}, dan 1.000 \textit{test}. Pipeline captioning menggunakan encoder CNN \textit{frozen} dan decoder RNN/LSTM.

\begin{table}[H]
    \centering
    \caption{Ringkasan Task Captioning}
    \begin{tabular}{ll}
        \toprule
        \textbf{Komponen} & \textbf{Spesifikasi} \\
        \midrule
        Task & Generasi caption gambar \\
        Arsitektur acuan & Show and Tell (\textit{pre-inject}) \\
        Decoder & SimpleRNN dan LSTM (dilatih terpisah) \\
        Loss & Sparse Categorical Crossentropy \\
        Optimizer (Keras) & Adam \\
        Metrik utama & BLEU-4 \\
        Metrik sekunder & METEOR, waktu eksekusi \\
        \bottomrule
    \end{tabular}
    \label{tab:ringkasan-caption}
\end{table}

\section{Batasan Implementasi}
Beberapa batasan yang digunakan agar konsisten dengan ketentuan tugas adalah sebagai berikut.

\begin{enumerate}
    \item Implementasi \textit{from scratch} untuk \textit{forward/backward} layer dilakukan menggunakan NumPy.
    \item Pelatihan model utama dilakukan di Keras, kemudian bobot dimuat ke implementasi \textit{scratch}.
    \item Encoder CNN pada tugas captioning digunakan dalam kondisi \textit{frozen}.
    \item Eksperimen dilakukan secara modular sehingga setiap variasi arsitektur/hyperparameter dapat dibandingkan secara adil.
\end{enumerate}
